% GENERATED FILE. Edit canonical lessons, then run npm run generate:books.
% canonical-source-hash: fnv1a64:5054c491bedede94
% canonical-lessons: ZH-C476-bozi, ZH-C476-shetou, ZH-C476-jianbang, ZH-C476-yaotou, ZH-C476-tengtong

\chapter{Neck and Shoulders}
\label{ch:zh-neck-and-shoulders}
\hlchaptermodality{476}

\begin{chapteropening}
\textbf{By the end of this chapter:} \emph{I can talk about my neck, tongue and shoulders, shaking my head, and pain getting better.}

Complete the last lesson of chapter 476: I can talk about my neck, tongue and shoulders, shaking my head, and pain getting better.
\end{chapteropening}

\section[bózi]{bózi — neck}

\label{lesson:ZH-C476-bozi}



\begin{quote}
Before the new one: say the Mandarin for a person's height, then the Mandarin for to bend.
\end{quote}

\subsection*{You'll want to know: bózi}
\emph{bózi} — \textquotedblleft{}neck\textquotedblright{}.

\begin{quote}
\emph{wǒ kàn le sān ge xiǎoshí diànnǎo bózi hěn téng} — I spent three hours at the computer, and my neck really hurts.
\end{quote}

Its characters come later; for now it is a word for the ear.

\begin{cousinweb}
\emph{bózi} is built from \emph{bó} and \zh{子}, a noun ending, both inherited from Middle Chinese, \emph{bwot} and \emph{tsiX} (Guangyun, in Baxter's transcription). The dictionary glosses the first as the navel; how it came to mean the neck is not settled.
\end{cousinweb}

\begin{grammarlens}[title={What you've built}]
One more word for saying how it was.
\end{grammarlens}

\subsection*{Your turn}
\begin{itemize}
\raggedright
  \item \emph{Say it:} \emph{bózi}
  \item \emph{Say it:} \emph{bózi}, once more
  \item \emph{Say it:} \emph{wǒ kàn le sān ge xiǎoshí diànnǎo bózi hěn téng} — I spent three hours at the computer, and my neck really hurts.
  \item \emph{Recall:} say the Mandarin for a person's height, then the Mandarin for to bend, then say \emph{bózi} again
\end{itemize}

\begin{tcolorbox}[breakable,colback=teal!4,colframe=teal!35!black,title={Before you move on}]
What does bózi mean? (\textquotedblleft{}Neck\textquotedblright{}.) Say it once more.
\end{tcolorbox}
\section[shétou]{shétou — tongue}

\label{lesson:ZH-C476-shetou}



\begin{quote}
Before the new one: say the Mandarin for to bend, then the Mandarin for the neck.
\end{quote}

\subsection*{You'll want to know: shétou}
\emph{shétou} — \textquotedblleft{}tongue\textquotedblright{}.

\begin{quote}
\emph{wǒ de shétou téng le liǎng tiān} — My tongue has been sore for two days.
\end{quote}

Its characters come later; for now it is a word for the ear.

\begin{cousinweb}
\emph{shétou} is built from \emph{shé} 'the tongue' and \zh{头}, a noun ending here, both inherited from Middle Chinese, \emph{zyet} and \emph{duw} (Guangyun, in Baxter's transcription). The dictionary lists the first under two readings and glosses \emph{zyet} as the tongue in the mouth; the second is said light.
\end{cousinweb}

\begin{grammarlens}[title={What you've built}]
One more word for saying how it was.
\end{grammarlens}

\subsection*{Your turn}
\begin{itemize}
\raggedright
  \item \emph{Say it:} \emph{shétou}
  \item \emph{Say it:} \emph{shétou}, once more
  \item \emph{Say it:} \emph{wǒ de shétou téng le liǎng tiān} — My tongue has been sore for two days.
  \item \emph{Recall:} say the Mandarin for to bend, then the Mandarin for the neck, then say \emph{shétou} again
\end{itemize}

\begin{tcolorbox}[breakable,colback=teal!4,colframe=teal!35!black,title={Before you move on}]
What does shétou mean? (\textquotedblleft{}Tongue\textquotedblright{}.) Say it once more.
\end{tcolorbox}
\section[jiānbǎng]{jiānbǎng — shoulder}

\label{lesson:ZH-C476-jianbang}



\begin{quote}
Before the new one: say the Mandarin for the neck, then the Mandarin for the tongue.
\end{quote}

\subsection*{You'll want to know: jiānbǎng}
\emph{jiānbǎng} — \textquotedblleft{}shoulder\textquotedblright{}.

\begin{quote}
\emph{wǒ de jiānbǎng hěn téng bù néng ná dōngxi} — My shoulder hurts so much I can't carry anything.
\end{quote}

Its characters come later; for now it is a word for the ear.

\begin{cousinweb}
\emph{jiānbǎng} is built from \emph{jiān} 'the shoulder' and \emph{bǎng}. The first is inherited from Middle Chinese \emph{ken} (Guangyun, in Baxter's transcription), glossed as the part below the nape. The dictionary lists the second only as \emph{bang}, in the word for the bladder, so the reading \emph{bǎng} and the sense 'shoulder' are not settled.
\end{cousinweb}

\begin{grammarlens}[title={What you've built}]
One more word for saying how it was.
\end{grammarlens}

\subsection*{Your turn}
\begin{itemize}
\raggedright
  \item \emph{Say it:} \emph{jiānbǎng}
  \item \emph{Say it:} \emph{jiānbǎng}, once more
  \item \emph{Say it:} \emph{wǒ de jiānbǎng hěn téng bù néng ná dōngxi} — My shoulder hurts so much I can't carry anything.
  \item \emph{Recall:} say the Mandarin for the neck, then the Mandarin for the tongue, then say \emph{jiānbǎng} again
\end{itemize}

\begin{tcolorbox}[breakable,colback=teal!4,colframe=teal!35!black,title={Before you move on}]
What does jiānbǎng mean? (\textquotedblleft{}Shoulder\textquotedblright{}.) Say it once more.
\end{tcolorbox}
\section[yáotóu]{yáotóu — to shake one's head}

\label{lesson:ZH-C476-yaotou}



\begin{quote}
Before the new one: say the Mandarin for the tongue, then the Mandarin for the shoulder.
\end{quote}

\subsection*{You'll want to know: yáotóu}
\emph{yáotóu} — \textquotedblleft{}to shake one's head\textquotedblright{}.

\begin{quote}
\emph{wǒ wèn tā qù bu qù tā yáotóu} — I asked him whether he was going, and he shook his head.
\end{quote}

Its characters come later; for now it is a word for the ear.

\begin{cousinweb}
\emph{yáotóu} is built from \emph{yáo} 'to shake' and \zh{头} 'head', both inherited from Middle Chinese, \emph{yew} and \emph{duw} (Guangyun, in Baxter's transcription), which glosses the first 'to move'.
\end{cousinweb}

\begin{grammarlens}[title={What you've built}]
One more word for saying how it was.
\end{grammarlens}

\subsection*{Your turn}
\begin{itemize}
\raggedright
  \item \emph{Say it:} \emph{yáotóu}
  \item \emph{Say it:} \emph{yáotóu}, once more
  \item \emph{Say it:} \emph{wǒ wèn tā qù bu qù tā yáotóu} — I asked him whether he was going, and he shook his head.
  \item \emph{Recall:} say the Mandarin for the tongue, then the Mandarin for the shoulder, then say \emph{yáotóu} again
\end{itemize}

\begin{tcolorbox}[breakable,colback=teal!4,colframe=teal!35!black,title={Before you move on}]
What does yáotóu mean? (\textquotedblleft{}To shake one's head\textquotedblright{}.) Say it once more.
\end{tcolorbox}
\section[téngtòng]{téngtòng — pain; ache}

\label{lesson:ZH-C476-tengtong}



\begin{quote}
Before the new one: say the Mandarin for the shoulder, then the Mandarin for to shake one's head.
\end{quote}

\subsection*{You'll want to know: téngtòng}
\emph{téngtòng} — \textquotedblleft{}pain; ache\textquotedblright{}.

\begin{quote}
\emph{chī yào yǐhòu téngtòng hǎo le yìdiǎnr} — After taking the medicine, the pain eased a little.
\end{quote}

Its characters come later; for now it is a word for the ear.

\begin{cousinweb}
\emph{téngtòng} is built from \zh{疼} 'pain' and \emph{tòng} 'to hurt', both inherited from Middle Chinese, \emph{downg} and \emph{thuwngH} (Guangyun, in Baxter's transcription), which glosses the first 'pain' and the second 'illness' and 'to hurt'.
\end{cousinweb}

\begin{grammarlens}[title={What you've built}]
One more word for saying how it was.
\end{grammarlens}

\subsection*{Your turn}
\begin{itemize}
\raggedright
  \item \emph{Say it:} \emph{téngtòng}
  \item \emph{Say it:} \emph{téngtòng}, once more
  \item \emph{Say it:} \emph{chī yào yǐhòu téngtòng hǎo le yìdiǎnr} — After taking the medicine, the pain eased a little.
  \item \emph{Recall:} say the Mandarin for the shoulder, then the Mandarin for to shake one's head, then say \emph{téngtòng} again
\end{itemize}

\begin{tcolorbox}[breakable,colback=teal!4,colframe=teal!35!black,title={Before you move on}]
What does téngtòng mean? (\textquotedblleft{}Pain; ache\textquotedblright{}.) Say it once more.
\end{tcolorbox}
